

We study the problem of fair $k$-committee selection under an egalitarian objective. Given $n$ agents partitioned into $m$ groups (\eg, demographic quotas), the goal is to aggregate their preferences to form a committee of size $k$ that guarantees minimum representation from each group while minimizing the maximum \emph{cost} incurred by any agent.
We model this setting as the ordinal fair $k$-center problem, where agents are embedded in an unknown metric space, and each agent reports a complete preference ranking (\ie, ordinal information) over all agents, consistent with the underlying distance metric (\ie, cardinal information).
The cost incurred by an agent with respect to a committee is defined as its distance to the closest committee member.
The quality of an algorithm is evaluated using the notion of distortion, which measures the worst-case ratio between the the cost of the committee produced by the algorithm and  the cost of an optimal committee, when given complete access to the underlying metric space. 


When cardinal information is not available, no constant distortion is possible for the ordinal $k$-center problem, even without fairness constraints,  when $k\geq 3$ [Burkhardt et.al., AAAI'24]. 
To overcome this hardness, we allow limited access to cardinal information by querying the metric space. In this setting, our main contribution is a factor-$5$ distortion algorithm that requires only $\bigO(k \log^2 k)$ queries. Along the way, we present an improved factor-$3$ distortion algorithm using 
$\bigO(k^2)$ 
queries.
